\section{Despliegue privado identificado por versión}\label{sec:impl-despliegue}

El proceso de entrega utiliza etiquetas de tres componentes numéricos con el
prefijo \texttt{v}. La automatización valida su formato y resuelve la etiqueta
al commit de origen antes de reutilizar las comprobaciones de Rust y de la
interfaz. Solamente después de su aprobación compila y empaqueta el ejecutable,
los archivos estáticos y la biblioteca qpdf fijada. Un manifiesto conserva
versión, commit, huellas de archivos y huella de migraciones; no contiene
credenciales ni material criptográfico privado.

La transferencia usa una identidad Ed25519 y una llave de servidor verificada
previamente. PostgreSQL, Redis, Rust y nginx se ejecutan como servicios del
usuario de despliegue, con puertos limitados a la interfaz de retorno local.
El navegador accede por un túnel SSH. Los datos, la clave de protección y el
material de firma y sellado permanecen fuera de los directorios de versiones.
La preparación del host no crea automáticamente una cuenta Owner ni publica
la aplicación en Internet.

La activación verifica el paquete, detiene la entrada y la API, respalda el
estado y cambia un enlace a la versión candidata. Después comprueba las
dependencias, la API, el proxy, el HTML y la identidad servida. Si falla,
restablece la versión anterior y conserva las escrituras de la base actual.
Una diferencia entre huellas de migraciones impide el cambio automático:
retroceder el ejecutable no implica que una base migrada sea compatible.
\rutarepo{docs/deployment.md} documenta mantenimiento, consulta de registros y
recuperación; \rutarepo{docs/adr/0048-private-versioned-deployment.md} fundamenta
esta separación entre versiones de aplicación y estado persistente.
